% ID: SEL-D0070-OB-E03
% ESTADO: cotejada. Segunda pasada visual y borrador compilado revisados.
% Fuente: Selectividad/sel_2009_matematicas_aplicadas/mat apli CC SS04.pdf
% SHA-256: bdea2fa323085e0db4eacff922ba9f403f58ef4e6bb4f2b569a0b857554c9ef2
% Página: 2; opción B; ejercicio original 3.

\noindent\underline{Parte I}

Un polideportivo dispone de 100 bolas de pádel y 120 bolas de tenis. Se sabe que 65 bolas son nuevas. Además, 75 bolas de pádel son usadas. Por un error, todas las bolas se han mezclado.
\begin{enumerate}[label=\alph*)]
\item \textbf{(1 punto)} Calcule la probabilidad de que si elegimos, al azar, una bola de tenis, ésta sea usada.
\item \textbf{(1 punto)} Calcule la probabilidad de que si elegimos, al azar, una bola, sea nueva.
\end{enumerate}
\par\noindent\begin{minipage}{\linewidth}
\noindent\underline{Parte II}
\begin{enumerate}[label=\alph*)]
\item \textbf{(1 punto)} En una población, una variable aleatoria $X$ sigue una distribución Normal de media 50 y desviación típica 9. Se elige, al azar, una muestra de tamaño 64 de esa población. \textquestiondown{}Cuál es la probabilidad de que la media muestral esté comprendida entre 48 y 52?
\item \textbf{(1 punto)} En una empresa de gas trabajan 150 personas en mantenimiento, 450 en operaciones, 200 en servicios y 100 en cargos directivos. Con objeto de realizar una encuesta laboral, se quiere seleccionar una muestra de 180 trabajadores de esa empresa por muestreo aleatorio estratificado con afijación proporcional, \textquestiondown{}qué número de trabajadores se debe elegir de cada grupo?
\end{enumerate}
\end{minipage}\par
